\subsection{Case 2608.20180: theorem structure and explicit families}\label{case:contrib-20180}

This mathematical-physics slice compares a general cyclic-geometry construction, its proof-level identities, and explicit solution families. Five decisions are recorded.

The off-shell Einstein--Hilbert density appears in the main argument and in a long appendix derivation. These are occurrences of one candidate, not two contributions. An occurrence overview may display one statement with both locators, reducing two display units to one without losing the appendix proof role.

The imported target-geodesic lifting principle and the paper's current-work geodesic sufficiency result overlap mathematically. They remain separate because attribution and proof role differ. Likewise, the general solution theorem, the Ricci-system identity, and the Tangherlini family are connected by implication and specialization, but none is a duplicate. The theorem packages quantified hypotheses; the identity supplies a load-bearing derivation; the family exhibits a concrete geometry and its interpretation.

This case rejects node-count novelty. A transitive mathematical skeleton could hide proof exposition only after a separate complete derivability witness; the generated fixture has no such claim-level hide. A contribution profile must distinguish the general theorem, its new proof mechanism, and explicit specialized constructions. The same verified closure can support different contribution operations.